% !TeX root = ../guide
\chapter{\University\ Thesis Requirements}\label{ch:requirements}
	\notice{%
		I do recognize that there are other types of theses that are accepted at the \University, however, this document will only be focusing on the Chapter- and Paper-Based Theses.}
	The following Sections will review and comment on the current requirements based on the various documents provided by the \Faculty \cite{FGPS2024,FGPS2025a,FGPS2024a,FGPS2025}.

	\section{Title Page}
		The \Faculty\ provides effectively zero wiggle room on the title page, requiring it strictly follow the sample they provided \cite{FGPS2025}.

	\section{Abstract}
		The abstract only has a few requirements:
		\begin{itemize}
			\item it must immediately follow the title page.
			\item it must be labelled as page ``ii''.
			\item it must be double spaced.
			\item it must \textbf{not} include graphs, charts, tables, or illustrations.
		\end{itemize}

	\section{Preface}
		\warning{%
			The inclusion of a preface does not excuse a student from failing to acknowledge the contributions of others in the body of the thesis, as per the University's Research and Scholarship Integrity Policy and the Code of Student Behaviour. 
			One would still expect to see footnotes, endnotes, or in-text references within the thesis acknowledging the works.}
		
		\caution{%
			Acknowledgements, such as thanks to the supervisor and supervisory committee members, to colleagues, lab mates and friends, and to family, \textbf{do not} appear in the preface.}
		
		\notice{%
			If you need assistance on writing the preface, ask your supervisor.Your supervisor must review and verify the preface before it becomes part of the final version of the thesis.}

		The requirements for the preface \cite{FGPS2025a}:
		\begin{itemize}
			\item Provide \textbf{full bibliographic citations} for any material in the thesis that has already been published and indicate the Chapter(s) where it appears.
			\item Acknowledge co-authorship or research collaboration that produced any of the material in the thesis and indicate the relative contributions of all collaborators and co-authors.
			\item Provide details of any ethics approval you obtained for the research.
			\item Acknowledge the use of generative Artificial Intelligence/Large Language Models if this has been permitted by your program and supervisory committee.
			\item Acknowledge third-party funding (grants, foreign governments, \textit{etc.}) or, 
			\item State that the thesis consists solely of your own unpublished and unfunded work.
		\end{itemize}

		The following sections will provide you with examples of how to include specific elements into your preface.
		The examples in the following sections have been modified from \fullcite{FGPS2024}.
		Due to the \enquote{Appendix for Thesis Formatting Guidelines} being generated using ChatGPT3.5 there are some inconsistencies within the document and the examples provided; while these are fine for providing a general idea of what to include, this goes against the ideology of this document.
		To combat this I have taken the liberty of adapting the examples to better blend into the style of this document and to increase the consistency, replacing some the generated examples with real examples where possible.

		Full credit still goes to the \University's Faculty of Graduate and Postdoctoral Studies for openly providing the document and I would strongly suggest anyone writing their thesis at the \University\ to read the updated documentation that they have put out in 2024.
		This new documentation provides much more guidance than the previous iteration which is no longer available.

		Before diving right into the examples in \Cref{preface:REA,preface:CW,preface:PPW,preface:AI,preface:CF,preface:DNA}, we will explore some of the still applicable guidelines for including your preface.
		First of which is a small change: since the FGSR became the \Fac, the requirement of a preface has become \textbf{mandatory} regardless of the contents of the thesis.

		When a thesis contains journal articles, the preface serves as a place for the student to include a statement indicating his or her contribution to the journal articles. 
		If any of the work presented in the thesis has led to any publications (accepted or published), these publications must be listed clearly in the preface with their bibliographical details and an indication as to where in the thesis this work is located (\textit{e.g.} state in which Chapter or Chapters). 
		For jointly authored publications\footnote{Permission may be needed if the co-authors hold the copyright in these publications.}, indication must also be given as to the relative contributions of the collaborators and co-authors, and a statement as to the proportion of research and writing conducted by the student.
		If ethics approval was required for the research, a statement to this effect must be included in the preface with the details of the approval that was granted. 

		\note{%
			Before preceding I would like to emphasize that the following examples are are not a one-and-done, you might need to combine elements from multiple of the following sections to build an appropriate preface for your thesis.}

		\info{%
			If \textbf{none} of the following apply \textbf{make sure} to include a preface in the style of the one shown in \Cref{preface:DNA}.}

		\tip{%
			Fun fact before you read on, Professor C.\ Ayranci who was mentioned in a few of the examples provided by \citetitle{FGPS2024} was in-fact the supervisor for my Master's thesis.}

		\subsection{Research Ethics Approval}\label{preface:REA}
			If your thesis required you to get Research Ethics Approval, then you should include a preface based on the one here.
	
			\begin{quote}
				\singlespacing
				\enquote{%
					This thesis is an original work by Daniel R.\ Aldrich. 
					The research project, of which this thesis is a part of, received research ethics approval from the \University\ Research Ethics Board 3, Project Name "Etymologies and Entomologies: Unraveling the Threads of Language and Ecology", No. 12345, February 23, 1993.}%TODO Change This to a fun project Name.
			\end{quote}

		\subsection{Collaborative Work}\label{preface:CW}
			If your thesis required you to work collaboratively with other, organizations, researchers or otherwise, then you should include a preface based on the one here.
	
			\begin{quote}
				\singlespacing
				\enquote{%
					Some of the research conducted for this thesis forms part of an international research collaboration, led by Professor T.\ Raivio at the University of Hogwarts, with Professor S.\ Agrawal being the lead collaborator at the \University. 
					The technical apparatus referred to in \Cref{ch:documentstructure} was designed by myself, with the assistance of Professor A.\ Shiri and Professor C.\ Ayranci. 
					The data analysis in \Cref{ch:figureandtables} and concluding analysis in \Cref{ch:plotsandgraphs} are my own work, as well as the literature review in \Cref{ch:gettingstarted}.}
			\end{quote}

		\subsection{Previously Published Material}\label{preface:PPW}
			If your thesis is based on or includes work previously published by yourself, or work you co-authored, then you should include a preface based on the one here.
	
			\begin{quote}
				\singlespacing
				\enquote{%
					\Cref{ch:gettingstarted} of this thesis has been published as \fullcite{Aldrich2017}. 
					I was responsible for the programming, data collection, and analysis as well as the manuscript composition. 
					C.\ Ayranci assisted with the data collection by providing access to the Multi-Functional Composites Group's Lab and Equipment, and contributed to manuscript edits. 
					D.\ Nobes was the supervisory and corresponding author and was involved with concept formation and manuscript composition.}
			\end{quote}

		\subsection{Use of Artificial Intelligence (AI)}\label{preface:AI}
			If your thesis used AI to help you outline, or otherwise write sections of your thesis; including help with analyzing, summarizing, \textit{etc.}, then you should include a preface based on the one here.
	
			\begin{quote}
				\singlespacing
				\enquote{%
					The generative artificial intelligence application or Large Language Model ChatGPT 3.5 was used for data analysis, summarization, synthesis, and simulation in \Cref{ch:documentstructure} of this thesis, as well as to generate a preliminary draft of the literature review in \Cref{ch:introduction}.}
			\end{quote}

		\subsection{Receiving of Competitive Funding}\label{preface:CF}
			If you received competitive funding used for your thesis, education, or research then you should include a preface based on the one here.
	
			\begin{quote}
				\singlespacing
				\enquote{%
					This work was supported by a Doctoral Fellowship from the Social Sciences and Humanities Research Council, a grant from the Entomological Association of Edmonton, and the National Scholarship Council of Narnia.}
			\end{quote}

		\subsection{Previous Prefaces Did Not Apply}\label{preface:DNA}
			\caution{%
				This was previously an \textit{optional} inclusion in your thesis, however, since the FGSR transitioned to the \Fac\ and have released new guidelines\cite{FGPS2025a} for formatting your thesis, this has now become \textbf{mandatory}.}

			If your thesis did not fall into any of the categories shown previously, then you should include a preface based on the one here. 

			\begin{quote}
				\singlespacing
				\enquote{%
					This thesis is an original work by \texttt{YOUR FULL NAME}. 
					No part of this thesis has been previously published.}
			\end{quote}

	\section{Dedications or Quotes}
		\notice{%
			Dedications or Quotations are limited to maximum of one page combined.}
		

	\section{Acknowledgements}
		\notice{%
			Acknowledgements are limited to maximum of two pages total.}

		The Acknowledgements is a strongly recommended, but not a strictly mandatory section of the thesis.
		For a thesis at the \University, an Acknowledgement page must be no more than two (2) pages in length.

		The Acknowledgements page serves as a place within a thesis where students may wish to acknowledge the provision of funding from third parties, such as an external scholarship bodies, research granting agencies, and foreign governments. 
		It is also appropriate to recognize the assistance provided by the supervisor and members of the supervisory committee.
		Additionally, you can recognize any help or support that you have received from family, friends, classmates, lab-mates, or any other individual that was able to help you in any way.
		A clear example of this would be the inclusion of the following, assuming that this template, class, or guide helped you in any significant way\footnote{Please be aware that it is not necessary to include this statement in any way, shape, or form it is only provided as an example on how to include these type of statements; this template/guide is provided in good faith with the \textbf{only} expectation being that you \textbf{will} share it with your fellow colleagues, supervisor, and or anyone else that may benefit from this.}:

		\begin{quote}
			\enquote{%
				I would like to thank Daniel R.\ Aldrich for his continuing contributions to the \University, and for his work within the Graduate Student Community. 
				More specifically, I would like to acknowledge the work that he has put into creating the \LaTeX\ template that this thesis was created in and the ongoing support that he provides to the students at the \University\ in formatting their theses.}
		\end{quote}

	\section{Table of Contents}
		Your Table of Contents must include all Chapter headings, 2--4 levels of subheadings, as well as all prefatory pages, excluding the title page.

	\section{List of Figures, Tables, Plates, \textit{etc.}}
		Include a list for each kind of non-textual item included in your thesis (figures, tables, plates, illustrations, videos, sounds files, \textit{etc.}).
		Each list must start on its own page.
		Each list must be numbered sequentially, and be in a consistent format (\textit{i.e.}, 1,2,3,... or 1.1,1.2,1.3,...).

	\section{List of Symbols or Abbreviations}
		Ask your supervisor or program director if your thesis needs a list of symbols or abbreviations. 
		Each list must start on a separate page.

	\section{Glossary of Terms}
		Ask your supervisor or program director if this is necessary.

	\section{Other Preliminary Items}
		Any discipline- or thesis-specific prefatory materials not mentioned in these guidelines should be placed directly before the body of the thesis.

	\section{Body of the Thesis}
		Minimum requirement: introduction, presentation of your research, and conclusion. The first page is always numbered with Arabic numeral 1. All subsequent pages are numbered consecutively.

	\section{Bibliography, Works Cited, or References}
		The Bibliography\footnote{Note: you may choose to have references at the end of each chapter or included as part of a paper-based thesis. In these cases, there may be some duplication of the works cited depending on how they are included.} has a few rules:
		\begin{itemize}
			\item it must be at the end of your thesis but before the appendices.
			\item it must \textbf{not} have a chapter number on it.
		\end{itemize}

	\section{Appendices}
		Supporting material that is referenced in the body of the thesis. 
		Appendices are not numbered as chapters, and should be kept to a minimum. 
		Pages are numbered consecutively with the Bibliography.